\section{Preliminaries}\label{sec:prelim}

\subsection{Magmas and equational laws}

\begin{definition}\label{def:magma}
A \emph{magma} is a pair $(M, \op)$ where $M$ is a nonempty set and
$\op \colon M \times M \to M$ is a binary operation.  No axioms are imposed on~$\op$.
\end{definition}

We write $x \op y$ for the application of the binary operation.  Terms over the
magma signature are built from variables $x, y, z, w, \ldots$ and the binary
operation~$\op$.  We denote the set of all such terms by $T(\op)$.

\begin{definition}\label{def:eqlaw}
An \emph{equational law} is an expression $s = t$ where $s, t \in T(\op)$.
The \emph{order} of a law is the total number of occurrences of~$\op$ in $s$
and~$t$.
\end{definition}

The Equational Theories Project \cite{ETP2025} catalogued all $4{,}694$
equational laws of order at most~$4$, up to renaming of variables and symmetry
of the equality sign.  We denote these laws $E_1, E_2, \ldots, E_{4694}$.

\begin{definition}\label{def:implication}
Let $E_i$ and $E_j$ be equational laws.  We say $E_i$ \emph{implies} $E_j$,
written $E_i \models E_j$, if every magma satisfying $E_i$ also
satisfies~$E_j$.  We say $E_i$ \emph{finitely implies} $E_j$, written
$E_i \models_{\mathrm{fin}} E_j$, if every \emph{finite} magma satisfying $E_i$
also satisfies~$E_j$.
\end{definition}

By a classical result of Birkhoff, $E_i \models E_j$ if and only if $E_j$ can
be derived from $E_i$ by finitely many applications of substitution and the
rewrite rule given by~$E_i$.  The relation $\models$ is a preorder on the set of
equational laws, and we denote the induced equivalence relation by~$\equiv$.

\begin{definition}\label{def:equiv-class}
Two laws $E_i$ and $E_j$ are \emph{equivalent}, written $E_i \equiv E_j$, if
$E_i \models E_j$ and $E_j \models E_i$.
\end{definition}

The $4{,}694$ laws form $1{,}415$ equivalence classes under~$\equiv$.  The
largest class contains $1{,}496$ laws, all equivalent to the \emph{singleton
law} $E_2 \colon x = y$.

\subsection{Structural classification of equational laws}

We now introduce the central definitions of this paper: a classification of
equational laws by their syntactic structure.

\begin{definition}\label{def:form}
Let $E \colon s = t$ be an equational law.  The \emph{structural form} of $E$
is one of the following five types:
\begin{enumerate}[label=(\roman*)]
\item \textbf{Trivial:} $s$ and $t$ are the same variable (e.g., $x = x$).
  There is exactly one such law, $E_1$.
\item \textbf{Singleton:} $s$ and $t$ are distinct variables (e.g., $x = y$).
  There is exactly one such law, $E_2$.
\item \textbf{Absorbing:} one side is a variable that does not appear on the
  other side (e.g., $x = y \op (z \op w)$).  There are $815$ such laws.
\item \textbf{Standard (lhs\_var):} one side is a variable that \emph{does}
  appear on the other side (e.g., $x = x \op (y \op z)$).  There are $2{,}322$
  such laws.
\item \textbf{General:} both sides contain at least one occurrence of~$\op$
  (e.g., $x \op y = y \op x$).  There are $1{,}555$ such laws.
\end{enumerate}
We write $\form(E) \in \{\text{triv}, \text{sing}, \text{abs}, \text{std},
\text{gen}\}$ for the structural form of~$E$.
\end{definition}

\begin{remark}\label{rem:form-partition}
The five forms partition the $4{,}694$ laws:
$1 + 1 + 815 + 2{,}322 + 1{,}555 = 4{,}694$.  By convention, when a law has
one side that is a bare variable, we take that side to be the left-hand side.
\end{remark}

\begin{definition}\label{def:signature}
The \emph{signature} of an equational law $E \colon s = t$ is the pair
$\sig(E) = (a, b)$ where $a$ is the number of occurrences of~$\op$ in~$s$ and
$b$ is the number of occurrences of~$\op$ in~$t$.
\end{definition}

\begin{example}\label{ex:signatures}
The law of associativity, $x \op (y \op z) = (x \op y) \op z$, has signature
$(2, 2)$.  The law $x = x \op (y \op z)$ has signature $(0, 2)$.  The
singleton law $x = y$ has signature $(0, 0)$.
\end{example}

\begin{definition}\label{def:nvar}
For an equational law $E$, we write $\nvar(E)$ for the number of
\emph{distinct variables} appearing in~$E$.
\end{definition}

\begin{example}\label{ex:nvar}
The law $x = x \op y$ has $\nvar = 2$.  The law
$x = (y \op z) \op (y \op w)$ has $\nvar = 4$.  The singleton law $x = y$
has $\nvar = 2$.
\end{example}

\subsection{Duality}

\begin{definition}\label{def:duality}
The \emph{dual} (or \emph{conjugate}) of a term $t \in T(\op)$ is obtained by
replacing every subterm $u \op v$ with $v \op u$.  The dual of a law
$E \colon s = t$ is the law $E^* \colon s^* = t^*$.
\end{definition}

Duality is a symmetry of the implication preorder: $E_i \models E_j$ if and
only if $E_i^* \models E_j^*$ \cite{ETP2025}.  This effectively halves the
number of implications that require independent verification.

\subsection{The implication matrix}

We denote the full implication matrix by $\mI \in \{0, 1\}^{4694 \times 4694}$,
where $\mI_{ij} = 1$ if and only if $E_i \models E_j$.  The following
statistics, computed from the complete matrix \cite{ETP2025}, inform our
analysis.

\begin{center}
\begin{tabular}{@{}lr@{}}
\toprule
Statistic & Value \\
\midrule
Total non-reflexive pairs & $22{,}028{,}942$ \\
True implications & $8{,}178{,}279$ (37.1\%) \\
False implications & $13{,}855{,}357$ (62.9\%) \\
Equivalence classes & $1{,}415$ \\
Laws equivalent to $E_2$ & $1{,}496$ \\
\bottomrule
\end{tabular}
\end{center}
